\section{Zoo Bridge Post-Quantum Attestations}
\label{sec:bridge-pq}

The \Zoo{} Bridge moves assets between \Zoo{} L2 and external
chains: \Lux{} C-Chain (canonical), Ethereum, BSC, and other EVM
ecosystems for which the bridge committee has stood up adapters.
In 2.0 the bridge attestation was a classical multisig over the
mint event. In 3.0 the bridge attestation lives on M-Chain and is
secured by post-quantum threshold signatures per
LP-016 \cite{lp-016} and LP-017 \cite{lp-017}.

\subsection{Bridge attestation flow under 3.0}

\begin{enumerate}[leftmargin=1.4em]
  \item \textbf{Source observation.} The bridge committee observes a
        finalised burn or lock event on the source chain.
  \item \textbf{Quorum vote.} The committee runs a threshold
        $\MLDSA{}\text{-}65$ ceremony on M-Chain: each member submits
        an $\MLDSA{}$ signature over the event digest. The threshold
        is $t$-of-$n$ with $t \geq \lceil 2n/3 \rceil$.
  \item \textbf{\Corona{} co-attestation.} The committee additionally
        emits a \Corona{} threshold proof over the same digest.
        Both proofs must succeed for a mint message to be valid.
  \item \textbf{B-Chain emission.} The pair of proofs is wrapped into
        a B-Chain bridge message (\textsc{lp-016}) and delivered
        across to the destination chain.
  \item \textbf{Destination mint.} The destination contract verifies
        both proofs --- $\MLDSA{}$ via on-chain precompile and
        \Corona{} via the verifier shipped in 3.0 --- and only
        then mints.
\end{enumerate}

\subsection{Why two PQ proofs on the bridge}

Bridges are the highest-value targets in the ecosystem. A break of
either lattice scheme would compromise mint authority on the
destination chain. Requiring both \Corona{} (Ring-LWE) and
$\MLDSA{}$ (Module-LWE/SIS) means the attacker must break two
distinct lattice problems --- not just two parameter sets of the
same family --- to forge a mint. Combined with the source-chain
finality requirement and the threshold $t \geq \lceil 2n/3 \rceil$
on the committee, this is the strongest practical defence against
both Q-day and committee compromise.

The classical $\BLS{}$ proof is also recorded for fast-path
verification on legacy explorers and as defence-in-depth, but
it is not load-bearing in 3.0; verifiers ignore it for mint
authorisation.

\subsection{Committee key schedule (LP-017)}

Per LP-017 \cite{lp-017}, the committee re-keys on a defined
schedule:

\begin{itemize}[leftmargin=1.2em]
  \item Every 90 days the committee runs a key-refresh ceremony,
        rotating each member's $\MLDSA{}$ and \Corona{} share.
  \item The ceremony is itself M-Chain-anchored: the new shares
        are committed under the previous round's threshold key,
        so an attacker who compromises a single member at one
        round cannot persist the compromise through the rotation.
  \item Members who fail to participate in two consecutive rotations
        are removed from the committee by L1 governance.
\end{itemize}

\subsection{Cross-chain message format}

A 3.0 bridge mint message carries:

\begin{lstlisting}[language=C++,caption={Bridge mint message v3}]
struct BridgeMintMsg3 {
    ChainId   src_chain;
    ChainId   dst_chain;
    Hash      src_event;            // finalised burn/lock event
    Hash      asset;
    Address   beneficiary;
    Amount    amount;
    Hash      committee_epoch;      // identifies key schedule epoch
    MlDsa::Sig    mldsa_threshold;  // t-of-n ML-DSA-65
    Corona::Sig corona_threshold; // t-of-n Corona
    Bls::Agg      bls_legacy;       // for fast-path observers, advisory
};
\end{lstlisting}

The destination-chain verifier rejects on a missing or invalid
$\MLDSA{}$ or \Corona{} proof. The $\BLS{}$ field is optional under
3.0.

\subsection{2021 BSC migration adapter}

The \Zoo{} Bridge supports a one-way migration adapter for legacy
2021-cohort BSC \Zoo{} holders. Holders burn the legacy BEP-20
token; the bridge committee attests the burn under the 3.0
attestation discipline; the holder receives the \Zoo{} 3.0 native
asset on \Zoo{} L2. The legacy-side burn cannot itself be PQ-attested
(BSC does not implement \MLDSA{} or \Corona{} verification), but
the destination-side mint is fully PQ-attested. This is the
correct posture: the chain we control is post-quantum; the chain
we are migrating from continues to be classical for as long as
upstream chooses.

\subsection{Failure modes and recovery}

\begin{itemize}[leftmargin=1.2em]
  \item \textbf{Quorum failure.} If the committee cannot reach the
        threshold, the message is not emitted. Mint does not happen.
        No silent fallback to a smaller quorum.
  \item \textbf{Single-member compromise.} The threshold structure
        absorbs up to $n - t$ compromised members. Compromise above
        that triggers L1 governance to slash and rotate.
  \item \textbf{Lattice scheme break.} If a published break of either
        \Corona{} or $\MLDSA{}$ surfaces, the bridge halts pending
        emergency re-parameterisation; the committee proposes new
        parameters via L1 governance and re-keys before resuming.
\end{itemize}
